% Pursuing Stacks (cotes 134-2 to 134-8): archive cote 134-4, site page 80 (PDF page 81), typescript page 168.
% typescript. Draft compared: ps3.tex:4097-4165.
% AI-transcribed by Claude Opus 5.5 (claude-opus-5-5, adaptive thinking, effort high), Message Batches API, batch ps2 (msgbatch_01NVwr1Yuo8qJf5rykeGBXPK), 2026-11-12; not yet checked by eye.
\scanpage{134-4/80}
\mspage{168}

To come to an end of this long terminological and notational digres-%
sion,  I'll have to say a word still about test categories,  modelizers,
test functors.  We surely want to use freely the terminology introduced
so far,  while we were working with ordinary weak equivalences,  in the
more general setting when a basic localizer $\underline{W}$ is given beforehand.  As
long as there is only one $\underline{W}$ around,  which will be used systematically in
all our constructions,  we'll just use the \inserted{previous} terminology,  being understood
that a ``modelizer'' say will mean a ``$\underline{W}$-modelizer'' ,  namely a category
M endowed with a \inserted{saturated} $W\unsure{\subset} \mathrm{F}\unsure{\mathrm{l}}(M))$,  such that $W^{-1}M$ is equivalent some way or
other to $\underline{W}^{-1}(\mathrm{Cat})$.  T\over{H}{h}e latter category however I dare not just designate
as (Hot),  as this notation has been associated to the very specific
situation of just ordinary homotopy types,  therefore I'll \inserted{always} write
$(\mathrm{Hot}_{\underline{W}})$  instead,  as a reminder of $\underline{W}$ after all !  If at a later moment
it should turn out that we'll have to work with more than one $\underline{W}$  \hand{(} for
instance,  to compare the $\underline{W}$-theory to the ordinary $W_{(\mathrm{Cat})}$ -theory),  we
will of course have to be careful and reintroduce $\underline{W}$ \xout{to specify} in our
wording,  to \inserted[typed]{qualify}\xout{specify} all notions dependant on the choice of $\underline{W}$ .

\hand{7.4.}

\hand{65}.  It has been over a week now and about eighty pages typing,since I
realized the need for looking at more general test functors than before
and hit upon how to handle them, that I am grinding stubbornly through
generalities unending on homotopy notions.  The grinding is a way of mine
to become familiar with a substance,  and at the \inserted{same} time getting aloof of it\unsure{\hand{,}}
climbing up,  sweatingly maybe,  to earn a birds eye view of a landscape
and maybe,  who knows,  in the end start a-flying in it, wholly at ease..
I am not there yet !  The least however one should expect\xout{\illegible{1}},  is that
the test story should now go through very smoothly. As I have been loo-%
sing contact lately with test categories and test functors,  I feel it'
ll be worth while to make still another review \inserted[typed]{of}\xout{upon} these notions,  lea-%
ding up to the key result I have been after all that time.  It will be
a way both to gain perspective,  and check if the grinding has been effi-%
cient indeed ...

If
\[ \xout{M}\ u\colon M \longrightarrow M' \]
is a functor between categories endowed each with a saturated set of
arrows,  W and W' say,  we'll say u is ``\uline{model preserv\xout{\illegible{1}}ing}''  (with respect
to (W,W')) if it satisfies the conditions:

\margin{80}
